\documentclass[10pt,a4paper]{amsart}
\usepackage[margin=19mm]{geometry}
\usepackage{amsmath,amssymb,mathtools}
\usepackage[scr=rsfs,cal=euler]{mathalfa}
\usepackage{xfrac}
\usepackage[unicode,hidelinks,bookmarks=false]{hyperref}
\newcommand{\ie}{i.e.\ }
\newcommand{\cf}{cf.\ }
\newcommand{\resp}{respectively\ }
\newcommand{\Subsection}{\subsection}
\providecommand{\nobreakdash}{\nobreak\hbox{-}}
\setlength{\emergencystretch}{2em}
\multlinegap0pt
\usepackage{amssymb}
\usepackage{float}
\usepackage[normalem]{ulem}
\newtheorem{lemma}{Lemma}
\title{Tree-like is not a transitive relation on paths}
\author{Jeremy Brazas, Gregory R. Conner, Paul Fabel, Curtis Kent}
\date{Source: arXiv:2603.16774v1 (CC BY 4.0)}
\begin{document}
\maketitle

\noindent\emph{Source and licence.} Tree-like is not a transitive relation on paths. Version arXiv:2603.16774v1 (CC BY 4.0), distributed by arXiv under the Creative Commons Attribution 4.0 International licence, http://creativecommons.org/licenses/by/4.0/, as recorded in the arXiv licence metadata for this exact version and shown on the abstract page https://arxiv.org/abs/2603.16774v1. Full licence notice: \texttt{licenses/arxiv-2603.16774-CC-BY-4.0.txt}.\par\smallskip

\noindent\textit{Editorial scope.} Counted unit: the article's lemma homotopic (source0.tex lines 366--379). The article's complete original proof of it is delivered with it (source0.tex lines 382--399, one \texttt{proof} environment of 18 lines). No mathematics is written, repaired or re-derived: the statement and the proof are the article's own lines, in the article's own notation, and no retained line is substituted or re-flowed.  Retained uncounted as the source-local context the proof needs: lines 331--333.  Nothing was omitted from any retained span, so the package carries no omission record.  No retained line contains a citation, so the wrapper carries no bibliography and no bibliography entry.  No retained line contains a figure environment, an image-inclusion command or drawing code, so no image is delivered and the package declares no asset dependency.  Editorial formatting only, in addition: an AMS \texttt{amsart} preamble that restates the article's own declarations and macros used by the retained lines, namely the environments \texttt{lemma} on separate counters, together with the standard packages those lines need.  The retained environments print in this document as Lemma 1 in order of appearance, whereas the article prints the same items with the numbering of its own full document.  The complete native source of the article as distributed in the arXiv e-print for this exact version is bundled unchanged as \texttt{sources/196577/source0.tex}.
\begin{lemma}\label{lem: the dendrite}
Suppose that $E_i$ is a nested sequence of dendrites with geodesic metrics $d_i$ such that $d_{i}|_{E_{i-1}\times E_{i-1}} = d_{i-1}$ and suppose $a>0$. If, for all $i>1$ , $\rho_i:E_i\to E_{i-1}$ is a monotone retract and $\sup\limits_{e\in E_i}d_i\big(\rho_i(e), e\big)\leq \frac{a}{2^i}$, then $E = \varprojlim(E_i, \rho_i)$ is a dendrite and  $E$ admits a metric $d$ such that  $d|_{E_i\times E_i}= d_i$ (where we identify $E_i$ with its natural embedding into $E$) and the projections $\varrho_i: E\to E_i$ converge uniformly to the identity map on $E$.
\end{lemma}

\begin{lemma}\label{lem: limit is r-tree homotopic}
Let $E_i$ be a nested sequence of dendrites with geodesic metrics $d_j$ such that $d_{i}|_{E_{i-1}\times E_{i-1}} = d_{i-1}$ and suppose $a>0$. Suppose that we have maps $\pi_i:[0,1]\to E_i$ and $g_i: E_i\to \mathbb R^2$ for $i\geq 1$ and maps $\rho_i:E_i\to E_{i-1}$ for $i\geq 2$ with the following properties.

    \begin{enumerate}
        \item $\rho_i:E_i\to E_{i-1}$ is a monotone retract of dendrites
        \item $\sup\limits_{e\in E_i}d_i\big(\rho_i(e), e\big)\leq \frac{a}{2^i}$
        \item\label{fixed ends}$\pi_i(r) = \pi_j(r)$ for all $i,j$ and $r\in \{0,1\}$
        \item\label{cauchy} $\sup_t d_i\big(\pi_{i-1}(t), \pi_i(t)\big)\leq \frac{a}{2^{i-1}}$.
        \item $g_i: E_i\to \mathbb R^2$ is $L$-Lipschitz.
        \item $g_i|_{E_{i-1}} = g_{i-1}$
    \end{enumerate}

    For $E = \varprojlim(E_i, \rho_i)$,  the maps $\pi_i$ converge uniformly to a map $\pi: [0,1]\to E$ such that $\pi(r) = \pi_1(r)$ for $r\in \{0,1\}$.  The maps $g_i\circ\varrho_i$ converge uniformly to an $L$-Lipschitz map $g: E \to \mathbb R^2$ such that $g|_{E_i} = g_i$.  In particular, $g\circ\pi$ is $\mathbb R$-tree homotopic to $g_1\circ \pi_1$
\end{lemma}

\begin{proof}
    By Lemma \ref{lem: the dendrite}, $E = \varprojlim(E_i, \rho_i)$ admits a metric $d$ such that  $d|_{E_i\times E_i}= d_i$ (where we identify $E_i$ with its natural embedding in $E$) and the projections $\varrho_i: E\to E_i$ converge uniformly to the identity map on $E$. By (\ref{cauchy}), the maps $\pi_i$ form a Cauchy sequence (in the sup metric) of continuous maps $[0,1]\to E$.  Thus they converge uniformly to a map $\pi:[0,1]\to E$ and (\ref{fixed ends}) implies that $\pi(r)= \pi_1(r)$ for $r\in \{0,1\}$.

    Let $e= (e_i)\in E$.  Then for $k>j$, we have

    \begin{align*}
        d\big(g_k\circ\varrho_k(e), g_j\circ\varrho_j(e)\big) &=d\big(g_k(e_k), g_j(e_j)\big)\\
                    &=d\big(g_k(e_k), g_k(e_j)\big)\leq  \sum\limits_{i=j+1}^{k}d\big(g_k(e_i), g_k(e_{i-1})\big)\\
                    &\leq\sum\limits_{i=j+1}^{k}L\cdot d_k\big(e_i, e_{i-1})\leq\sum\limits_{i=j+1}^{k}L\frac{a}{2^i} \leq L\frac{a}{2^j}.
    \end{align*}


    Thus $g_k\circ\varrho_k$ is a Cauchy sequence of maps $E\to \mathbb{R}^2$ (in the sup metric) and therefore converge to a map $g: E\to \mathbb{R}^2$.  Notice that $g_k\circ\varrho_k|_{E_k} = g_k$ for all $k$.  Hence for all $\ell\geq k$, $g_\ell\circ\varrho_\ell|_{E_k} = g_k$.  Thus $g|_{E_k} = g_k$ which implies that $g\circ \pi_1 = g_1\circ \pi_1$.  Since $E$ is a dendrite, $\pi_1(0)= \pi(0)$, and $\pi_1(1) = \pi(1)$, the path $g\circ \pi_1$ is $\mathbb R$-tree homotopic to $g\circ \pi$. Hence $g_1\circ \pi_1$ is $\mathbb R$-tree homotopic to $g\circ \pi$.


    Since $\varrho_i$ is $1$-Lipschitz and $g_i$ is $L$-Lipschitz, $g_i\circ \varrho_i$ is $L$-Lipschitz which implies that $g$ is $L$-Lipschitz.

\end{proof}

\end{document}
